\documentclass[10pt,a4paper]{amsart}
\usepackage[margin=19mm]{geometry}
\usepackage{amsmath,amssymb,mathtools}
\usepackage[scr=rsfs,cal=euler]{mathalfa}
\usepackage{xfrac}
\usepackage[unicode,hidelinks,bookmarks=false]{hyperref}
\newcommand{\ie}{i.e.\ }
\newcommand{\cf}{cf.\ }
\newcommand{\resp}{respectively\ }
\newcommand{\Subsection}{\subsection}
\providecommand{\nobreakdash}{\nobreak\hbox{-}}
\setlength{\emergencystretch}{2em}
\multlinegap0pt
\usepackage{amssymb}
\usepackage{tikz}
\newtheorem{proposition}{Proposition}
\newcommand{\E}{{\mathbb E}}
\newcommand{\N}{{\mathbb N}}
\newcommand{\V}{{\mathcal M}}
\title{Beta distribution and associated Stirling numbers of the second kind}
\author{Jakub Gismatullin, Patrick Tardivel}
\date{Source: arXiv:2601.08453v1 (CC BY 4.0)}
\begin{document}
\maketitle

\noindent\emph{Source and licence.} Beta distribution and associated Stirling numbers of the second kind. Version arXiv:2601.08453v1 (CC BY 4.0), distributed by arXiv under the Creative Commons Attribution 4.0 International licence, http://creativecommons.org/licenses/by/4.0/, as recorded in the arXiv licence metadata for this exact version and shown on the abstract page https://arxiv.org/abs/2601.08453v1. Full licence notice: \texttt{licenses/arxiv-2601.08453-CC-BY-4.0.txt}.\par\smallskip

\noindent\textit{Editorial scope.} Counted unit: the article's proposition prop:bounds\_expo (source0.tex lines 171--184). The article's complete original proof of it is delivered with it (source0.tex lines 500--525, one \texttt{proof} environment of 26 lines, which the article places away from the statement). No mathematics is written, repaired or re-derived: the statement and the proof are the article's own lines, in the article's own notation, and no retained line is substituted or re-flowed.  No further source-local context is needed: every cross-reference of every retained line resolves inside the two delivered spans.  Nothing was omitted from any retained span, so the package carries no omission record.  No retained line contains a citation, so the wrapper carries no bibliography and no bibliography entry.  No retained line contains a figure environment, an image-inclusion command or drawing code, so no image is delivered and the package declares no asset dependency.  Editorial formatting only, in addition: an AMS \texttt{amsart} preamble that restates the article's own declarations and macros used by the retained lines, namely the environments \texttt{proposition} on separate counters, together with the standard packages those lines need.  The retained environments print in this document as Proposition 1 in order of appearance, whereas the article prints the same items with the numbering of its own full document.  The complete native source of the article as distributed in the arXiv e-print for this exact version is bundled unchanged as \texttt{sources/192595/source0.tex}.
\begin{proposition}
\label{prop:bounds_expo}
Let $k \in \N$, $m\in \N_{>0}$, $r\in \N_{>1}$ and $\mathcal{E}_1,\dots,\mathcal{E}_m$ be i.i.d random variables having standard exponential distribution with density $\exp(-x)$.
\begin{itemize}
    \item[i)] The following inequality holds
    \begin{equation}
    r^k\V_r(k,m)\le \left(\frac{r}{r-1}\right)^{2k} \E\left(\mathcal{E}_1+\dots+\mathcal{E}_m\right)^k =\left(\frac{r}{r-1}\right)^{2k} \frac{(m-1+k)!}{(m-1)!}.
    \end{equation}
    \item[ii)] The upper bound given in i) is sharp since the following limit holds 
    \begin{equation}
    \lim_{r\to +\infty} r^{k}\V_r(k,m)=\E\left(\mathcal{E}_1+\dots+\mathcal{E}_m\right)^k= \frac{(m-1+k)!}{(m-1)!}.
    \end{equation}
\end{itemize}
\end{proposition}

\begin{proof}{Proposition {\ref{prop:bounds_expo}}}
i) Note that the density of the random variable $Y_i=(r-1)X_i$  on $[0,r-1]$ is
\[ f(x) =\frac{r}{r-1}\left(1-\frac{x}{r-1}\right)^{r-1}.\]
Due to the following inequality for all $x\in [0,r-1]$ 
\[
\frac{r}{r-1}\left(1-\frac{x}{r-1}\right)^{r-1} \le \frac{r}{r-1}\exp(-x)
\]
one may deduce that 
\begin{eqnarray*} (r-1)^k\E\left(\sum_{i=1}^m X_i  \right)^k = \E\left(\sum_{i=1}^m Y_i  \right)^k &\le& \left(\frac{r}{r-1}\right)^k \E\left(\sum_{i=1}^m \mathcal{E}_i\right)^k\\
&\le& \left(\frac{r}{r-1}\right)^k \frac{(m-1+k)!}{(m-1)!}
\end{eqnarray*}since $\sum_{i=1}^m \mathcal{E}_i$ has an Erlang distribution  (the density of $\sum_{i=1}^m \mathcal{E}_i$ is  \\
$h(x)=\frac{x^{m-1}\exp(-x)}{(m-1)!}$) whose  moment of order $k$ is $\frac{(m-1+k)!}{(m-1)!}$. This inequality completes the proof of i).

ii) The density of the random variable $(r-1)X_1$ converges pointwise to the density  of $\mathcal{E}_1$ namely $\lim_{r\to +\infty}f(x)=\exp(-x)$, for all $x\in [0,+\infty)$. Therefore, the following limit holds
\[ 
  \lim_{r\to +\infty}x^kf^{*m}(x)=
  \frac{x^{m+k-1}\exp(-x)}{(m-1)!}.
\]
Finally, the Dominated Convergence Theorem gives that
\begin{eqnarray*}
\lim_{r\to +\infty}
(r-1)^k\E\left(\sum_{i=1}^m X_i\right)^k &=&
\lim_{r\to +\infty} \int_{0}^{+\infty} \frac{x^{m+k-1}\exp(-x)}{(m-1)!}dx\\
&=&  \E\left(\sum_{i=1}^m \mathcal{E}_i\right)^k=\frac{(m-1+k)!}{(m-1)!}.
\end{eqnarray*}
\end{proof}

\end{document}
